% Pursuing Stacks (cotes 134-2 to 134-8): archive cote 134-2, site page 67 (PDF page 68).
% appendix (letters to Breen). Draft compared: ps1.tex:1527-1544.
% AI-transcribed by Claude Opus 5.5 (claude-opus-5-5, adaptive thinking, effort high), Message Batches API, batch ps1 (msgbatch_01Rv5hCaNTtUt4JHx1yqoHPL), 2026-11-12; not yet checked by eye.
\scanpage{134-2/67}
\mspage{\unsure{52}}

\margin{\unsure{(App. 14)|(Apr. 14)}} (F) \uline{Significance and limitations of the fppf topology}

Since the attempts of Serre to find a ``Weil cohomology'' by using the
cohomology of a scheme with coefficients not only discrete $\mathbb{Z}/p^n\mathbb{Z}$ $(n \to \infty)$
\margin{$\mu_{p^n}$} or $\mu_{p^{\over{\ast}{n}}}$ $(n \to \infty)$ , but also continuous (for example $W_n$, $n \to \infty$) , which
give good results for recovering a correct $H^1$ , during numerous years I
have come upon the impression, which I have tried in vain to make precise,
that a \inserted{correct ``p-adic''} ``Weil cohomology''\hand{,} \deleted{``p-adic''}, in the case $p > 0$ \inserted{and k of char. p}, should come, in
one way or another, from the fppf cohomology, for finite coefficients
for example, or more general coefficients, e.g. algebraic groups over $k$ .
The construction in (B) of the local jacobian complex was,of course, related
to this hope: the homology might reveal what is hidden to us in cohomology!
For some time now, one has at one\hand{\unsure{'}}s disposal the formalism of crystalline
cohomology, and one knows (Berthelot) that (at least for $X$ projective
and smooth) it has the correct properties. If one uses that as a kind of
standard by which to ``measure'' the other cohomologies, one finds that
the part of the crystalline cohomology $H^i_{\text{cris}}(X)$ which could be described in
terms of fppf cohomology of $X$ with coefficients in \deleted{an} algebraic k-group\hand{s}.
\deleted{comes from} \inserted{is} a small part of $H^i$ \inserted{only} ; more precisely, using the very rigid
supplementary structure of the $H^i$ (modules of finite type on the ring $W(k)$
of Witt vectors) which comes from the existence of the Frobenius homomorphism \inserted{(an isogeny)}
$H^i \hand{\xrightarrow{F}} H^i$ (semi-linear)\hand{,} one finds that \deleted{they remain} \inserted{one keeps} always in the part ``of slope
$\leq 1$'' (although the possible slopes vary between 0 and i ...). This explains
why for $i = 1$ one can obtain via fppf a correct $H^1$ , although for $H^2$
already all the attempts have been unfruitful.

In truth, one conjectures that \textit{all} the part\deleted{s of points} \inserted{of slope} $\leq 1$ in $H^{\hand{i}}_{\text{cris}}$
come\hand{s} from fppf . But I have completely lost contact with these questions -
people such as Mazur, Katz, Messing - and of course Deligne - should be
knowledgeable as to the present states of these questions.

\begin{center}------\$------\end{center}

\margin{espacer}
\margin{\unsure{(App. 15)|(Apr. 15)}} Your question 7 seems to indicate that there is a misunderstanding on
your part on the significance of the ``homotopy type'' of $X$ , for $X$ a topos

\margin{67}
